\documentclass[11pt]{article}

\usepackage{amsmath,amssymb,amsthm,mathtools}
\usepackage[T1]{fontenc}
\usepackage{lmodern}
\usepackage[margin=1in]{geometry}
\usepackage{enumitem}
\usepackage{hyperref}
\usepackage{microtype}
\usepackage{booktabs}

\hypersetup{colorlinks=true,linkcolor=blue,citecolor=blue,urlcolor=blue}

\newtheorem{theorem}{Theorem}[section]
\newtheorem{proposition}[theorem]{Proposition}
\newtheorem{remark}[theorem]{Remark}

\newcommand{\Z}{\mathbb Z}

\title{D5 intended quotient identification after 105:\\
theorem-side quotient versus dynamic bridge}
\author{}
\date{March 2026}

\begin{document}
\maketitle

\begin{abstract}
This note records the right downstream statement of the D5 intended quotient
question after the odd-$m$ marked-chain descent note `105`.

The main point is negative before it is positive. The theorem-side intended
quotient should not yet be frozen as literal raw state-space equality with the
dynamic bridge $(\beta,q_0,\sigma)$ or $(\beta,\delta)$. The correct theorem
target is still the exact deterministic quotient on the exact marked object
that retains grouped base.

What the new check `106` does show is sharper and usable: on the regenerated
carry-jump-to-wrap state-chain object, the theorem-side key $(B,\beta)$
determines the dynamic bridge state $(\beta,q_0,\sigma)$ exactly, while the
dynamic bridge forgets source-residue bookkeeping and becomes equivalent to the
theorem-side key only after adjoining $\rho$.
\end{abstract}

\section{Executive answer}

The downstream intended quotient should be formulated as:
\[
Q_{\mathrm{int}}
=
\text{exact deterministic quotient on the exact marked object retaining grouped base.}
\]

The dynamic bridge
\[
Q_{\mathrm{dyn}}=(\beta,q_0,\sigma)\cong(\beta,\delta),\qquad \delta=q_0+m\sigma,
\]
is the right checked splice-compatible comparison target, but not yet the
literal global definition of $Q_{\mathrm{int}}$.

The regenerated comparison check on $m=13,15,17$ shows:
\begin{enumerate}[label=\textup{(\arabic*)},leftmargin=2.8em]
\item $(B,\beta)$ determines $(\beta,q_0,\sigma)$ exactly;
\item $(\beta,q_0,\sigma)$ is exact for current $\epsilon_4$, exact for the
      short-corner detector, and deterministic on successor;
\item $(\beta,q_0,\sigma)$ is not globally equal to $(B,\beta)$ because it
      forgets source-residue bookkeeping;
\item adjoining $\rho$ repairs that loss exactly:
      \[
      (B,\beta)\cong (\rho,\beta,q_0,\sigma)
      \]
      on the extracted exact marked chain object;
\item adjoining only the family bit does not repair the loss.
\end{enumerate}

So the right statement is not ``the intended quotient is exactly
$(\beta,\delta)$.'' The right statement is: the intended quotient factors to
the dynamic bridge, and after adjoining $\rho$ the two parameterizations are
equivalent on the exact marked object.

\section{Checked comparison package}

\begin{proposition}[checked comparison on regenerated chain objects]
On the regenerated carry-jump-to-wrap state-chain object for
\[
m\in\{13,15,17\},
\]
the following hold.
\begin{enumerate}[label=\textup{(\alph*)},leftmargin=2.5em]
\item The map
      \[
      (B,\beta)\longmapsto (\beta,q_0,\sigma)
      \]
      is exact.
\item The dynamic bridge $(\beta,q_0,\sigma)$ is exact for current
      $\epsilon_4$, exact for the short-corner detector, and deterministic on
      nonterminal successor.
\item Every dynamic bridge bucket is ambiguous as a theorem-side state.
\item The enriched bridge
      \[
      (\rho,\beta,q_0,\sigma)
      \]
      determines $(B,\beta)$ exactly.
\item The weaker enrichment
      \[
      (\mathrm{family},\beta,q_0,\sigma)
      \]
      does not determine $(B,\beta)$ exactly.
\end{enumerate}
\end{proposition}

\begin{proof}
This is the content of the saved comparison script
\texttt{torus\_nd\_d5\_intended\_quotient\_compare\_106.py} and its JSON
outputs in \texttt{RoundY/checks/d5\_106\_intended\_quotient\_compare/}.
\end{proof}

\begin{remark}
The ambiguity is systematic, not sporadic. For each checked modulus, the number
of ambiguous dynamic buckets is exactly $m^3$, and the bucket sizes lie in
\[
\{m-5,\ m-4,\ m-3\}.
\]
So the dynamic bridge is really missing a whole source-residue bookkeeping
layer.
\end{remark}

\section{What theorem should be proved next}

The right theorem program is now short.

\begin{enumerate}[label=\textup{(\arabic*)},leftmargin=2.8em]
\item Prove symbolically that on the exact marked object, the descended
      theorem-side data $(G,\beta)$ determines $(q_0,\sigma)$.
\item Prove that the only additional bookkeeping forgotten by the dynamic
      bridge is the source residue $\rho$.
\item State the comparison theorem in the form
      \[
      Q_{\mathrm{int}}\cong (\rho,\beta,q_0,\sigma)
      \quad\Longrightarrow\quad
      Q_{\mathrm{int}}\twoheadrightarrow (\beta,q_0,\sigma)\cong(\beta,\delta).
      \]
\end{enumerate}

This is the honest closure of full M3. It avoids the overstatement that the
intended quotient has already been proved to be literally $(\beta,\delta)$ as a
raw state space, while still capturing the exact checked relation that the
proof package now supports.

\section{Bottom line}

The two parameterizations are not globally identical as currently stated.
What is checked, and what now looks theorem-level plausible, is the sharper
comparison:
\[
(B,\beta)\cong(\rho,\beta,q_0,\sigma)
\quad\text{and hence}\quad
(B,\beta)\twoheadrightarrow(\beta,q_0,\sigma)\cong(\beta,\delta).
\]

So the next D5 quotient theorem should be a comparison theorem with $\rho$, not
a premature identity claim for bare $(\beta,\delta)$.

\end{document}
